Este informe propuso una matheurística que coordina la planificación semanal de recursos con el ruteo diario de una red de transporte sanitario, y la evaluó en las 80 semanas disponibles y en 92 corridas de sensibilidad y de infraestructura. Las conclusiones se ordenan según los objetivos específicos de la Sección~\ref{sec:objetivos}.

\parrafo{O1: coordinar ahorra cerca de un 10\,\%.} La matheurística es más barata que el caso base en las 80 semanas: un 9,9\,\% en BIG (\$3,9 millones por semana) y un 15,7\,\% en SMALL. El ahorro es estadísticamente significativo (intervalo de confianza de 95\,\% de 9,5 a 10,3\,\% en BIG) y no depende de la semilla ni de la utilización con que se dimensiona el caso base. Viene sobre todo de rutear mejor y no de contratar más: con casi la misma flota se recorre un 16\,\% menos de kilómetros y se externaliza un 28\,\% menos. Todas las soluciones pasan la validación independiente.

\parrafo{O2: el recurso escaso es el furgón frío.} En 26 de las 40 semanas BIG se usan los 18 furgones, y en esas semanas se externaliza más. Las solicitudes difíciles son casi siempre las mismas para ambos métodos: insumos y equipos grandes con ventanas de retiro de una hora, que necesitan un furgón o dos vehículos a la vez. La sensibilidad confirma el diagnóstico: más demanda, menos flota o externalización más cara encarecen la red justo cuando se agotan los furgones.

\parrafo{O3: la ventaja es robusta; el costo depende de la demanda y la congestión.} La matheurística es más barata que el caso base en todos los escenarios evaluados. El costo responde en forma convexa a la demanda ($+36$\,\% con 30\,\% más de solicitudes en la semana mediana) y sube hasta un 11\,\% con una congestión de 50\,\%, por lo que los costos reportados, calculados con flujo libre, son un piso. Ampliar las ventanas de atención en media hora baja el costo entre 3,5 y 4,6\,\% y casi elimina las externalizaciones; acortarlas en 15 minutos lo sube cerca de un 2,6\,\%. Entre las decisiones que la organización controla, exigir puntualidad total y partir la jornada a las 07:00 tienen costo nulo o negativo.

\parrafo{O4: la red actual está bien dimensionada, salvo en furgones.} Ninguna base es imprescindible (cerrar cualquiera sube el costo a lo más 1,2\,\%), ninguna de las tres ubicaciones evaluadas para una décima base paga su costo fijo y la capacidad de las bases no limita. La inversión que sí se justifica son dos furgones adicionales en las semanas de alta demanda, que bajan el costo un 1,2\,\% y las externalizaciones de 20 a 12 en la semana más exigente. Que cerrar algunas bases baje levemente el costo indica que la red podría operar con siete u ocho, pero la diferencia está dentro del margen de la heurística y debe confirmarse con más semanas.

\parrafo{O5: la división sostiene la operación y el método escala.} Sin división el costo sube un tercio, porque cerca de un quinto de las solicitudes no cabe en ningún vehículo. Los hubs, en cambio, casi no se usan aunque se ofrezcan transferencias para todas las solicitudes transferibles: con estos datos el desvío y la espera no compensan. Pasar de SMALL a BIG duplica las solicitudes y triplica la flota, pero el tiempo mediano solo sube de 16 a 25 minutos por semana, y la generación de rutas se puede repartir por día entre procesadores.

\parrafo{Fortalezas y limitaciones.} La herramienta resuelve una semana completa en unos 25 minutos en un computador personal, da el mismo resultado en cualquier equipo y entrega todo lo que la organización necesita para contratar y operar. Sus limitaciones son cuatro. La solución es aproximada y no hay una cota válida para el problema completo. Los tiempos de viaje suponen flujo libre. La base de cada vehículo la fija la heurística al generar las rutas, lo que deja ahorros sin capturar, como mostró el cierre de bases. Y el modelo supone que las solicitudes de la semana se conocen de antemano.

\parrafo{Pasos siguientes.} Se recomienda implementar la herramienta en las tres etapas del Capítulo~\ref{cap:implementacion}: una marcha blanca que mida los tiempos reales de viaje, un piloto en las bases del poniente evaluado con los cuatro criterios propuestos y el despliegue completo. En paralelo, las extensiones de mayor impacto son incorporar tiempos de viaje según la hora y dejar que el maestro elija dónde ubicar cada vehículo.
